\cleardoublepage % memastikan bab baru dimulai di halaman ganjil (kanan)
\chapter{PENDAHULUAN}
\label{chap:pendahuluan}

% --- Latar Belakang ---
\section{Latar Belakang}

%TODO-CITE: paragraf ini butuh 1-2 sitasi tentang tren elektrifikasi transportasi
% publik di Indonesia/dunia (mis. target elektrifikasi bus BRT, peta jalan
% kendaraan listrik nasional). Cari dari referensi yang sudah kamu kumpulkan
% dan tambahkan entry-nya ke daftar-pustaka.bib.
Elektrifikasi armada bus menjadi salah satu langkah yang mulai ditempuh
operator transportasi publik di Indonesia untuk mengurangi emisi dan biaya
operasional bahan bakar. TransJogja, sebagai operator Bus Rapid Transit di
Daerah Istimewa Yogyakarta, telah memulai tahap awal elektrifikasi ini dengan
mengoperasikan 2 unit bus listrik pada Rute L-1 (Terminal Jombor, Malioboro,
Terminal Jombor), mengisi daya di satu Stasiun Pengisian Kendaraan Listrik
Umum (SPKLU) di kawasan Adisucipto.

Dinas Perhubungan (Dishub) DIY mempertimbangkan perluasan elektrifikasi ini
ke seluruh 20 koridor\footnote{Dalam tugas akhir ini, istilah
\textit{koridor} dan \textit{rute} dipakai bergantian untuk merujuk pada
unit operasional yang sama -- salah satu dari 20 kode rute TransJogja
(mis. 1A, 1B, 2A, dst.) -- dan tidak membedakan tingkat hierarki
keduanya sebagaimana mungkin didefinisikan secara resmi oleh Dishub DIY.}
TransJogja, yang berarti sekitar 124 sampai 128 bus
(termasuk cadangan) akan beroperasi penuh bertenaga listrik. Perluasan pada
skala ini
membawa persoalan yang berbeda secara mendasar dari operasi 2 bus di 1 rute
saat ini. Bus listrik memiliki jangkauan terbatas per satu kali pengisian
penuh (untuk spesifikasi bus TransJogja, sekitar 150 km) dan waktu pengisian
yang jauh lebih lama dibandingkan pengisian bahan bakar bus diesel. Jika
lokasi dan kapasitas SPKLU tidak direncanakan dengan tepat, bus berisiko
kehabisan daya di tengah rute, mengantre lama untuk mengisi daya, atau
terpaksa memotong jumlah ritasenya, yang pada akhirnya menurunkan
keteraturan jadwal (headway) yang dirasakan penumpang.

Sampai saat penelitian ini dilakukan, belum ada studi kuantitatif yang
menjawab pertanyaan operasional mendasar untuk skenario elektrifikasi penuh
ini: pada titik mana di sepanjang rute bus harus mengisi daya, berapa banyak
SPKLU yang dibutuhkan dan di lokasi mana, berapa kapasitas (jumlah SPKLU)
tiap SPKLU harus disediakan, dan mekanisme pengisian daya apa (saat beroperasi
di sela rit, di depot, atau kombinasi keduanya) yang paling menjaga headway
tetap sesuai target perencanaan Dishub DIY. Pengujian langsung di lapangan
untuk menjawab pertanyaan-pertanyaan ini tidak praktis. Dibutuhkan investasi
infrastruktur SPKLU dan pengadaan bus dalam jumlah besar sebelum hasilnya bisa
diamati, dengan risiko kerugian operasional bila konfigurasi yang dipilih
ternyata keliru.

%TODO-CITE: 1 sitasi tentang penggunaan simulasi/Discrete Event Simulation
% untuk perencanaan infrastruktur transportasi, sebagai justifikasi kenapa
% pendekatan simulasi dipilih dibanding pengujian langsung.
Simulasi menawarkan jalan keluar dari keterbatasan ini. Dengan memodelkan
perilaku drain baterai tiap bus per segmen rute, mekanisme antrean di SPKLU,
dan variasi kondisi lalu lintas harian, berbagai skenario lokasi, kapasitas,
dan mekanisme pengisian daya dapat diuji secara cepat dan berulang tanpa
biaya maupun risiko dari pengujian lapangan. Tugas akhir ini menyusun
simulasi tersebut: model drain baterai non-linear berbasis literatur,
mesin simulasi berbasis peristiwa diskrit (\textit{Discrete Event Simulation})
untuk seluruh 20 koridor TransJogja, dan lapisan stokastik Monte Carlo untuk
menangkap variasi waktu tempuh harian yang tidak bisa ditangkap oleh simulasi
deterministik satu skenario saja.

% --- Rumusan Masalah ---
\section{Rumusan Masalah}

Perluasan elektrifikasi TransJogja dari 2 bus di 1 koridor menjadi seluruh
20 koridor membawa persoalan lokasi, kapasitas, dan mekanisme pengisian daya
yang belum terjawab secara kuantitatif. Ketiadaan jawaban ini berisiko
membuat Dishub DIY mengambil keputusan investasi infrastruktur SPKLU
berdasarkan asumsi kualitatif semata. Misalnya, menempatkan SPKLU hanya di
terminal-terminal utama yang sudah ada tanpa memperhitungkan pola permintaan
pengisian daya riil per koridor, yang berujung pada kegagalan menjaga
headway operasional setelah elektrifikasi penuh berjalan, atau pada
pembangunan infrastruktur yang berlebih (dan karenanya boros) di lokasi yang
sebenarnya tidak menjadi titik kebutuhan utama.

Berdasarkan uraian tersebut, rumusan masalah pada tugas akhir ini adalah
sebagai berikut:
\begin{enumerate}
	\item Bagaimana memodelkan drain baterai non-linear untuk bus listrik
	TransJogja yang tervalidasi secara kualitatif terhadap data operasional
	riil Rute L-1?
	\item Bagaimana membangun mesin simulasi operasi armada yang dapat
	menjalankan seluruh 20 koridor TransJogja secara simultan, termasuk
	mekanisme antrean di SPKLU?
	\item Bagaimana pengaruh skenario lokasi SPKLU terhadap headway operasional
	dan tingkat kegagalan rute pada kondisi elektrifikasi penuh?
	\item Berapa kapasitas SPKLU (jumlah SPKLU per lokasi) yang dibutuhkan tiap
	lokasi, dan pada titik mana penambahan atau pengurangan kapasitas mulai
	berdampak signifikan terhadap keandalan operasional?
	\item Bagaimana pengaruh mekanisme pengisian daya (\textit{opportunity
		charging}, \textit{depot charging}, dan kombinasi keduanya) terhadap
	headway operasional dan kebutuhan penambahan armada?
	\item Bagaimana hasil simulasi dapat disajikan secara visual agar mudah
	dibaca dan dikomunikasikan kepada pemangku kepentingan?
\end{enumerate}

% --- Tujuan ---
\section{Tujuan}

Tujuan utama tugas akhir ini adalah menghasilkan rekomendasi kuantitatif
mengenai lokasi, kapasitas, dan mekanisme pengisian daya SPKLU untuk skenario
elektrifikasi penuh 20 koridor TransJogja, melalui simulasi drain baterai dan
operasi armada. Secara lebih rinci, tugas akhir ini bertujuan untuk:
\begin{enumerate}
	\item Membangun model drain baterai non-linear untuk bus listrik TransJogja
	yang divalidasi secara kualitatif\footnote{Validasi kualitatif di sini
	berarti membandingkan kesesuaian arah dan orde besaran hasil simulasi
	terhadap data riil, bukan kecocokan numerik presisi -- lihat metode
	lengkap pada Bab V.} terhadap data operasional riil Rute L-1.
	\item Membangun mesin simulasi operasi armada berbasis peristiwa diskrit
	yang dapat menjalankan seluruh 20 koridor TransJogja secara simultan,
	termasuk mekanisme antrean di SPKLU.
	\item Membandingkan skenario lokasi SPKLU (kondisi saat ini, rencana yang
	sudah ada, serta kandidat lokasi baru berbasis analisis data permintaan)
	terhadap dampaknya pada headway dan tingkat kegagalan rute.
	\item Membandingkan skenario kapasitas SPKLU (jumlah SPKLU per lokasi) untuk
	menemukan titik ketika penambahan atau pengurangan kapasitas mulai
	berdampak signifikan terhadap keandalan operasional (titik patah/\textit{regime
		shift}).
	\item Membandingkan skenario mekanisme pengisian daya (\textit{opportunity
		charging}, \textit{depot charging}, serta \textit{mix charging} sebagai
	kombinasi keduanya) pada cakupan evaluasi harian (\textit{single-day},
	lihat Batasan Masalah), beserta kebutuhan penambahan armada yang menyertainya.
	\item Menyediakan \textit{dashboard} visual yang menampilkan hasil simulasi dalam bentuk
	peta skematik rute dan animasi pergerakan bus per skenario, sebagai alat
	bantu komunikasi hasil kepada pemangku kepentingan.
\end{enumerate}

Kriteria keberhasilan tugas akhir ini adalah tersedianya perbandingan
kuantitatif (headway aktual terhadap target Dishub DIY, jumlah rute yang
mengalami kegagalan pengisian daya, dan kebutuhan kapasitas SPKLU) untuk
minimal 3 skenario lokasi dan 4 skenario mekanisme pengisian daya, disertai
model drain baterai yang tervalidasi secara kualitatif terhadap baseline
Rute L-1.

% --- Batasan Masalah ---
\section{Batasan Masalah}

Untuk menjaga tugas akhir ini tetap dapat diselesaikan dalam waktu yang
tersedia, beberapa batasan berikut diterapkan:
\begin{enumerate}
	\item Simulasi tidak memodelkan keputusan pengemudi atau dispatch secara
	\textit{real-time} di lapangan; pergerakan bus mengikuti jadwal rit yang
	dibangkitkan dari data rute.
	\item Skenario lokasi SPKLU merupakan \textit{input} yang dibandingkan satu
	sama lain (kondisi saat ini, rencana yang sudah ada, dan kandidat baru
	berbasis analisis data), bukan hasil optimasi otomatis dari algoritma
	\textit{greenfield} (algoritma yang mencari lokasi optimal dari nol,
	tanpa asumsi infrastruktur eksisting). Pemilihan lokasi kandidat dilakukan melalui
	perbandingan skenario (\textit{grid search}/\textit{scenario comparison}),
	bukan algoritma metaheuristik atau \textit{reinforcement learning} formal.
	\item Simulasi tidak memodelkan efek cuaca atau perilaku penumpang secara
	mendetail; variasi kondisi lalu lintas direpresentasikan sebagai parameter
	variabilitas waktu tempuh.
	\item Tugas akhir ini tidak mencakup studi kelayakan finansial atas investasi
	infrastruktur SPKLU maupun pengadaan armada.
	\item Simulasi bersifat harian dan berdiri sendiri per hari (\textit{single-day/stateless}):
	status daya baterai tiap bus selalu diinisialisasi penuh pada awal hari
	simulasi, tanpa membawa sisa hasil pengisian daya dari malam sebelumnya.
	Akibatnya, perbandingan mekanisme pengisian daya pada tugas akhir ini
	efektif menguji dua rezim operasional siang hari, yaitu \textit{opportunity
		charging} tersedia versus tidak tersedia, dan bukan mengukur ketiga
	mekanisme (\textit{opportunity}, \textit{depot}, \textit{mix}) sebagai tiga
	kondisi yang sepenuhnya independen lintas hari.
\end{enumerate}

% --- Metodologi Pengerjaan TA ---
\section{Metodologi}

Tugas akhir ini dikerjakan mengikuti kerangka \textit{Design Thinking}, yang
terdiri atas lima tahap: \textit{Empathize}, \textit{Define}, \textit{Ideate},
\textit{Prototype}, dan \textit{Test}. Kelima tahap ini bersifat iteratif.
Temuan pada tahap \textit{Test} beberapa kali membawa proses kembali ke tahap
\textit{Ideate} untuk merevisi rancangan skenario sebelum diuji ulang. Kelima
tahap tersebut diadaptasi ke konteks pengerjaan tugas akhir sebagai berikut:
\begin{enumerate}
	\item \textbf{Empathize (Pemahaman Kondisi Saat Ini):} investigasi kondisi
	operasional TransJogja saat ini melalui data rute (jarak antar-halte,
	jumlah bus, target headway per koridor), spesifikasi bus listrik, dan
	aturan operasional pengisian daya, yang diperoleh dari dokumen internal
	Dishub DIY/operator dan hasil wawancara dengan pihak terkait. Studi
	literatur mengenai model konsumsi energi kendaraan listrik dan metode
	simulasi peristiwa diskrit untuk sistem transportasi turut dilakukan pada
	tahap ini sebagai dasar pemilihan pendekatan. Hasil tahap ini dijelaskan
	pada Bab II.
	\item \textbf{Define (Perumusan Masalah):} menganalisis data dan hasil studi
	literatur dari tahap \textit{Empathize} untuk merumuskan persoalan secara
	spesifik, menentukan baseline validasi (Rute L-1), dan menyusun daftar
	skenario lokasi, kapasitas, serta mekanisme pengisian daya yang akan diuji.
	Hasil tahap ini dijelaskan pada Bab III.
	\item \textbf{Ideate (Perancangan Solusi):} merancang model drain baterai,
	arsitektur mesin simulasi berbasis peristiwa diskrit, struktur data
	skenario, dan rancangan dasbor visual, berdasarkan hasil perumusan masalah.
	Hasil tahap ini dijelaskan pada Bab IV.
	\item \textbf{Prototype (Implementasi):} membangun purwarupa model dan mesin
	simulasi, mengalibrasi model drain baterai terhadap baseline Rute L-1, dan
	mengimplementasikan dasbor visual berbasis web. Hasil tahap ini dijelaskan
	pada Bab V.
	\item \textbf{Test (Evaluasi):} menjalankan simulasi Monte Carlo (repetisi
	hari dengan variasi waktu tempuh acak) untuk tiap skenario lokasi,
	kapasitas, dan mekanisme pengisian daya, dibandingkan berdasarkan headway
	aktual, tingkat kegagalan rute, dan metrik antrean SPKLU. Hasil pengujian
	yang menunjukkan kekurangan pada rancangan skenario tertentu ditindaklanjuti
	dengan kembali ke tahap \textit{Ideate} untuk perbaikan. Hasil tahap ini
	dijelaskan pada Bab VI.
\end{enumerate}
Kesimpulan dan saran yang ditarik dari keseluruhan proses ini dijelaskan pada
Bab VII.

% --- Sistematika Penulisan ---
\section{Sistematika Penulisan}

Laporan tugas akhir ini disusun dalam tujuh bab dengan sistematika sebagai
berikut:
\begin{enumerate}
	\item Bab I Pendahuluan: berisi latar belakang, rumusan masalah, tujuan,
	batasan masalah, metodologi, dan sistematika penulisan.
	\item Bab II Studi Literatur: berisi kajian mengenai model konsumsi energi
	kendaraan listrik, metode simulasi peristiwa diskrit, dan penelitian
	terkait perencanaan infrastruktur pengisian daya kendaraan listrik.
	\item Bab III Analisis: berisi analisis kebutuhan data, penentuan baseline
	validasi, dan rancangan skenario pengujian lokasi, kapasitas, serta
	mekanisme pengisian daya SPKLU.
	\item Bab IV Perancangan: berisi rancangan model drain baterai, arsitektur
	mesin simulasi, struktur data skenario, dan rancangan dasbor visual.
	\item Bab V Implementasi: berisi penjelasan implementasi model, mesin
	simulasi, kalibrasi terhadap baseline Rute L-1, dan dasbor visual.
	\item Bab VI Evaluasi: berisi hasil simulasi Monte Carlo dan perbandingan
	skenario lokasi, kapasitas, dan mekanisme pengisian daya SPKLU terhadap
	headway operasional dan tingkat kegagalan rute.
	\item Bab VII Kesimpulan dan Saran: berisi kesimpulan dari hasil tugas akhir
	dan saran untuk pengembangan lebih lanjut.
\end{enumerate}